\section{Mostowski Collapse}

In \Cref{Ch1:Thm:Mostowski_WO} we saw that every well-ordering is isomorphic to a unique ordinal, by a map that sends each point to the set of images of its predecessors. In this section we prove the general version promised there: the ordering hypothesis is dropped, and ``ordinal'' becomes ``transitive class''. One more property of the relation is needed to make the map injective, and we begin with it.

\subsection{Extensional Relations}

\begin{boxdefinition}[Extensionality]
    $(A, R)$ is \textbf{extensional} iff $\forall y, z \in A$,
    \begin{align*}
        \setst{x \in A}{x \, R \, y} = \setst{x \in A}{x \, R \, z}
        \to y = z
    \end{align*}
\end{boxdefinition}

This is just like the axiom of extensionality, but worded as a property of sets and relations.

Not all relations are extensional.

\begin{boxnexample}[Non-Extensional Relations] \hfill
    If $A = \set{\set{1}, \set{2}}$, then $(A, \in)$ is not extensional. The reason is that
    \begin{align*}
        \setst{x \in A}{x \in \set{1}} = \emptyset = \setst{x \in A}{x \in \set{2}}
    \end{align*}
    but obviously $\set{1} \ne \set{2}$.
\end{boxnexample}

But many are.

\begin{boxexample}[Extensional Relations] \hfill
    \begin{enumerate}
        \item If $M$ is a transitive class, then $(M, \in)$ is extensional. Indeed, suppose $y, z \in M$ and $y \ne z$. Since $M$ is transitive, $y, z \ssq M$, so $y = y \cap M$ and $z = z \cap M$. Thus
            \begin{align*}
                y \cap M = \setst{x \in M}{x \in y} \neq \setst{x \in M}{x \in z} = z \cap M
            \end{align*}

        \item If $(A, R)$ is a class linear ordering (ie, a linear ordering on the class $A$), then $(A, R)$ is extensional.

        For instance, if $A \ssq \OR$ and $R$ is the usual ordering of ordinals (write $R$ as $\in$ or $<$), then $(A, R)$ is extensional.
    \end{enumerate}
\end{boxexample}

\subsection{The Collapse Theorem}

We are now ready to state and prove Mostowski Collapse. Note that it will use neither Choice nor Foundation.

\begin{boxtheorem}[Mostowski Collapse Theorem Scheme]\label{Ch2:Thm:Mostowski}
    Let $(A, R)$ be a class structure that's well-founded, set-like and extensional. Define, using the theorem scheme on definitions by recursion, the mapping $G : A \to V$ by
    \begin{align*}
        \forall y \in A, G(y) = \setst{G(x)}{x \in A \text{ and } x \, R \, y}
    \end{align*}
    Write the range of $G$ as $M$. Then,
    \begin{enumerate}
        \item $M$ is a transitive class
        \item $G$ is a class isomorphism from $(A, R)$ to $(M, \in)$
        % [CORRECTED] was: $G' : (A, R) \simeq (M, \in)$ --- with $M$ rather than $M'$ the hypothesis never mentions $M'$ and the conclusion $M' = M$ fails for any other transitive class; the proof below uses that $M'$ is the range of $G'$
        \item If $M'$ is a transitive class and $G' : (A, R) \simeq (M', \in)$, then $M' = M$ and $G' = G$.
    \end{enumerate}
    Note that $G$ is referred to as the \textbf{Mostowski Collapse Isomorphism}.
\end{boxtheorem}
\begin{remark}
    When we say that $M' = M$ and $G' = G$ above, we are not saying that their defining formulae are the same. In other words, their equality is not definitional.
\end{remark}
\begin{proof}
    We may assume $R \ssq A \times A$.

    \begin{enumerate}
        \item Say $a \in b \in M$. We must see $a \in M$. Then, $b \in M$. As $M$ is the range of $G$, there's some $y \in  A$ such that $b = G(y)$. As $G(y) = \setst{G(x)}{x \, R \, y}$, there's some $x \, R \, y$ such that $a = G(x)$. Now $a = G(x) \in M$.

        \item We need to show $G$ is an order-preserving bijection.
        
        \begin{description}
            \item[\underline{Bijectivity.}] It's clear that $G$ is surjective, because $M$ is the range of $G$. To show $G$ is injective, it is enough to show that for every $x \in  A$,
            \begin{align}
                \forall x' \in A \brac{G(x) = G\of{x'} \to x = x'}
                \label{Ch2:Eq:Mostowski-pf-expr}
            \end{align}
            Let's prove \eqref{Ch2:Eq:Mostowski-pf-expr} by induction on $R$. Let $y \in A$ and assume \eqref{Ch2:Eq:Mostowski-pf-expr} holds for all $x \, R \, y$. We'll now prove that \eqref{Ch2:Eq:Mostowski-pf-expr} holds for $y$ as well.
    
            Fix any $y' \in A$ and assume that $G(y) = G\of{y'}$. We need to prove $y = y'$. Because $(A, R)$ is extensional, it is enough to show that
            \begin{align*}
                \setst{x}{x \, R \, y} = \setst{x'}{x' \, R \, y'}
            \end{align*}
            \begin{description}
                \item[$(\subseteq)$]
                    Suppose $x \, R \, y$. Then
                    \begin{align*}
                        G(x) \in G(y) = G\of{y'} = \setst{G\of{x'}}{x' \, R \, y'}
                    \end{align*}
                    So there's some $x' \, R \, y'$ with $G(x) = G\of{x'}$. By \eqref{Ch2:Eq:Mostowski-pf-expr} for $x$, we know that $x = x'$. Thus $x \, R \, y'$.
    
                \item[$(\supseteq)$]
                    Suppose $x' \, R \, y'$. Then,
                    \begin{align*}
                        G\of{x'} \in G\of{y'} = G(y) = \setst{G(x)}{x \, R \, y}
                    \end{align*}
                    So there's some $x \, R \, y$ such that $G\of{x'} = G(x)$. By \eqref{Ch2:Eq:Mostowski-pf-expr} for $x$, $x' \, R \, y$.
            \end{description}

            \item[\underline{Monotonicity.}] Suppose $G(x) \in G(y)$. By definition of $G(y)$, there is some $x' \, R \, y$ such that $G(x) = G\of{x'}$. As $G$ is injective, we know that $x = x'$. Thus, $x \, R \, y$.
        \end{description}

        \item Suppose $M'$ is a transitive class and $G' : (A, R) \simeq (M', \in)$. We need $M = M'$ and $G = G'$.

        Since $M$ is the range of $G$ and $M'$ is the range of $G'$, it is enough to show that $G = G'$. By the uniqueness of recursive definitions (\Cref{Ch2:Thm:Recursion}), it is enough to see that $G'$ satisfies the same defining equation as $G$. That is, $G'(y) = \setst{G'\of{x}}{x \, R \, y}$.

        \begin{description}
            \item[$(\supseteq)$]
                $x \, R \, y \to G'(x) \in G'(y)$ just because $G'$ preserves order (by assumption).

            \item[$(\subseteq)$]
                Suppose $a \in G'(y)$. Now $G'(y) \in M'$ trivially because $M'$ is the range of $G'$. Since $M'$ is a transitive class, $G'(y) \ssq M'$. Hence, $a \in M'$. So we have some $x \in A$ such that $a = G'(x)$. Now $G'(x) \in G'(y)$. Since $G'$ is an isomorphism, this implies $x \, R \, y$.
        \end{description}
    \end{enumerate}
\end{proof}

\subsection{Examples}

Let's study some examples. The first collapses a proper class, and what it collapses onto is a class we have had since \Cref{Ch1:CH}.

\begin{boxexample}[Sequences of Ordinals]
    Consider the following classes:
    \begin{align*}
        ^{< \omega}\OR &= \setst{f}{\exists n < \omega, \, f : n \to \OR} \\
        \OR^{< \omega} &= \setst{t}{\exists n < \omega, \, t \text{ is an $n$-tuple of ordinals}} \\
        \brac{\OR}^{< \omega} &= \setst{s}{\exists n < \omega , \abs{s} = n \text{ and } s \subset \OR}
    \end{align*}
    The above are pairwise in bijection. No parameters are necessary for anything here.

    % [CORRECTED] the clause fixing the length convention for $<_{\lex}$ is inserted; the text stated none. The other reading (a proper initial segment counts as \textit{larger}) makes $\emptyset$ the largest element, so the order type is $\OR + 1$ and the collapse cannot land on $(\OR, <)$, which has no largest element. Set-likeness survives that reading everywhere except at $\emptyset$, whose predecessors are then a proper class
    % [CORRECTED] was: ``where if $s_0 = \emptyset$'' --- a fragment with no consequent, and the only claim in it is false: $s_0$ is $\max(s)$, so $s_0 = \emptyset$ says $s = \set{0}$, which is the case $n = 1$. What the identification needs at this point is the degenerate case $n = 0$, where the displayed list has no terms and $s$ is the empty set
    Let's consider $\brac{\OR}^{< \omega}$. Say $s \in \brac{\OR}^{< \omega}$. Say $n < \omega$ and $s \in \brac{\OR}^{n}$. List $s$ in decreasing order as $s_0 > s_1 > \cdots > s_{n - 1}$, where $s = \emptyset$ if $n = 0$. Identify $s$ with the $n$-tuple $\parenth{s_0, \ldots, s_{n-1}}$. Let $<_{\lex}$ denote the lexicographic ordering on members of $\brac{\OR}^{< \omega}$ under this identification, where a tuple counts as smaller than any tuple properly extending it. We see that
    \begin{itemize}
        \item $<_{\lex}$ is a class well-ordering of $\brac{\OR}^{< \omega}$ that is definable without parameters
        \item $\parenth{\brac{\OR}^{<\omega}, <_{\lex}}$ is well-founded and extensional
        \item $<_{\lex}$ is set-like: given $s \in \brac{\OR}^{< \omega}$,
        \begin{align*}
            \setst{t \in \brac{\OR}^{< \omega}}{t <_{\lex} s} \ssq \brac{\max(s) + 1}^{< \omega}
        \end{align*}
    \end{itemize}
    We can apply the Mostowski Collapse Theorem Scheme (\Cref{Ch2:Thm:Mostowski}) to get a map
    \begin{align*}
        G : \parenth{\brac{\OR}^{< \omega}, <_{\lex}} \simeq \parenth{\OR, <}
    \end{align*}
    In particular, $G$ is a class bijection between $\brac{\OR}^{< \omega}$ and $\OR$ that is definable without parameters.
\end{boxexample}

The second example is less a computation than a technique, and it is the one we will keep coming back to.

\begin{boxexample}[ZFC]
    Consider any $\alpha \in \OR$. Recall that $V_{\alpha}$ is a transitive set, so $\parenth{V_{\alpha}, \in}$ is extensional, set-like and well-founded.

    Now, you can formalize all of mathematics in ZFC, including Model Theory---in particular, the downwards Löwenheim-Skolem Theorem (DLS). (We will do this eventually in this course.)

    Indeed, within the ZFC formalization, we can define the notion of an \textbf{elementary substructure} as follows: we say $X \preceq Y$ iff for every ``formula''\footnote{There are no formulae because we're in a formalism where everything's a set! But there are ``formulae'': just make air quotes and change your tone of voice and everything is ok :D} $\varphi\of{\bar{u}}$ and every matching $\bar{a}$ from $X$, $X \models \varphi\of{\bar{a}}$ iff $Y \models \varphi\of{\bar{a}}$. You can imagine that with similar handwaving, we can formalize other notions, like DLS.

    Once we have done all of this, we can prove that there is some countable $X$ such that $(X, \in) \preceq \parenth{V_{\alpha}, \in}$. Moreover, note that
    \begin{itemize}
        \item if $\alpha \leq \omega$ then $X = V_{\alpha}$
        \item if $\alpha \geq \omega + 1$, then $X \subsetneq V_{\alpha}$ ($X$ is countable but $V_{\alpha}$ is uncountable).
    \end{itemize}
    Observe that $(X, \in)$ is
    \begin{description}
        \item[\underline{well-founded:}] $X \subseteq V_{\alpha} \to X \in V_{\alpha + 1} \subset \WF$

        \item[\underline{set-like:}] clear

        \item[\underline{extensional:}] we know that $V_{\alpha}$ is transitive, hence extensional. This is equivalent to the following fact (\textit{within the ZFC-formalism of ZFC}):
        \begin{align*}
            V_{\alpha} \models \text{``Extensionality Axiom''}
        \end{align*}
        Then, since $X \preceq V_{\alpha}$,
        \begin{align*}
            X \models \text{``Extensionality Axiom''}
        \end{align*}
        as well. This is equivalent to $(X, \in)$ being extensional.

        We can also argue more directly for extensionality, which asks that any two distinct members of $X$ be told apart by a member of $X$. Say $b, c \in X$ with $b \ne c$. Since $V_{\alpha}$ is extensional, some $a \in V_{\alpha}$ lies in the symmetric difference $b \symdiff c$, so that
        \begin{align*}
            V_{\alpha} \models \exists a \parenth{a \in b \symdiff c}
        \end{align*}
        This is a ``formula'' whose parameters $b$ and $c$ both come from $X$, and $X \preceq V_{\alpha}$, so $X$ satisfies it too. The witness it supplies is a member of $X$ lying in $b \symdiff c$, which is exactly what we wanted.
    \end{description}
    Let $G : \parenth{X, \in} \simeq \parenth{M, \in}$ be the Mostowski Collapse Isomorphism (ie, if $G$ is the Mostowski Collapse Isomorphism, then $(M, \in)$ is the Mostowski Collapse---the range of $G$). As this is an isomorphism, $M$ is transitive. Note that $G\inv : M \to V_{\alpha}$ is an elementary embedding (we compose $G\inv : M \simeq X$ with the inclusion $X \inj V_{\alpha}$). In particular, $M$ and $V_{\alpha}$ satisfy the same ``sentences''.

    Intuitively, what's going on is that we have built some countable microcosm---some tiny, countable analog---of the entire $V$-universe, \textit{inside of the ZFC formalism}. This is one of the most important subroutines in set theory: the idea that you can take a big initial piece of $V$, take some elementary substructure via DLS, take its Mostowski Collapse, and do something with it. We will see a lot of this later on.

    \begin{figure}[H]
        \centering
        \begin{tikzpicture}[every node/.style = {font = \small}]
            \draw[thick] (2.6, 4.4) -- (0, 0) -- (-2.6, 4.4);
            \draw[thick, dashed] (-2.6, 4.4) -- node[above] {$V_{\alpha}$} (2.6, 4.4);
            \draw[thick, dotted] (-2.6, 4.4) -- (-3.1, 5.2);
            \draw[thick, dotted] (2.6, 4.4) -- (3.1, 5.2);
            \node at (0, 5.2) {$V$};

            \draw[thick] (-6.2, 3.2) -- (-6.9, 2.0) -- (-7.6, 3.2);
            \draw[thick, dashed] (-7.6, 3.2) -- node[above] {$M$} (-6.2, 3.2);

            \draw[->, thick] (-6.0, 2.6) -- node[above] {$G\inv$} (-1.6, 2.6);
        \end{tikzpicture}
    \end{figure}
\end{boxexample}
